\documentclass[12pt,a4paper]{extarticle}
\usepackage[utf8]{inputenc}
\usepackage[T5]{fontenc}
\usepackage{geometry}
\geometry{a4paper, left=1.5cm, right=1cm, top=1.5cm, bottom=1.5cm, headheight=20pt, headsep=15pt, footskip=28pt}
\usepackage{enumitem}
\usepackage{amssymb}
\usepackage{tikz}
\usetikzlibrary{patterns, patterns.meta} % Thêm patterns.meta vào đây
% Gọi package nằm trong thư mục con template
\usepackage{../template/mngocbox}

\begin{document}

% Sử dụng ngoặc vuông [...] để thay đổi linh hoạt các thông số
\makeheadbox[headerright={Tổng ôn 10-11},chude={CHỦ ĐỀ 04},tieude={Tam thức bậc hai},dinhdang={Đại số},lop={12 },gv={Nguyễn Lê Minh Ngọc}]

\vspace{2em}
\makeheading{IV}{Đề kiểm tra}
\noindent\textbf{Câu 1:} Tập nghiệm của bất phương trình $x^2-2x+1\ge 0$ là
\begin{center}
\begin{tabular}{p{3.5cm} p{3.5cm} p{3.5cm} p{3.5cm}}
\textbf{A.} $\varnothing$. & \textbf{B.} $\mathbb{R}$. & \textbf{C.} $(1;3)$. & \textbf{D.} $\mathbb{R} \setminus \{1\}$.
\end{tabular}
\end{center}
\noindent\textbf{Câu 2:} Tập nghiệm của bất phương trình $x^2-2x+1>0$ là
\begin{center}
\begin{tabular}{p{3.5cm} p{3.5cm} p{3.5cm} p{3.5cm}}
\textbf{A.} $\varnothing$. & \textbf{B.} $\mathbb{R}$. & \textbf{C.} $\{1\}$. & \textbf{D.} $\mathbb{R} \setminus \{1\}$.
\end{tabular}
\end{center}
\noindent\textbf{Câu 3:} Tập nghiệm của bất phương trình $x^2-2x+1<0$ là
\begin{center}
\begin{tabular}{p{3.5cm} p{3.5cm} p{3.5cm} p{3.5cm}}
\textbf{A.} $\varnothing$. & \textbf{B.} $\mathbb{R}$. & \textbf{C.} $\{1\}$. & \textbf{D.} $\mathbb{R} \setminus \{1\}$.
\end{tabular}
\end{center}
\noindent\textbf{Câu 4:} Tập nghiệm của bất phương trình $x^2-2x+1\le 0$ là
\begin{center}
\begin{tabular}{p{3.5cm} p{3.5cm} p{3.5cm} p{3.5cm}}
\textbf{A.} $\varnothing$. & \textbf{B.} $\mathbb{R}$. & \textbf{C.} $\{1\}$. & \textbf{D.} $\mathbb{R} \setminus \{1\}$.
\end{tabular}
\end{center}
\noindent\textbf{Câu 5:} Bất phương trình $x^2-3x-4\le 0$ có tất cả bao nhiêu nghiệm nguyên dương?
\begin{center}
\begin{tabular}{p{3.5cm} p{3.5cm} p{3.5cm} p{3.5cm}}
\textbf{A.} 6. & \textbf{B.} 5. & \textbf{C.} 4. & \textbf{D.} 2.
\end{tabular}
\end{center}
\noindent\textbf{Câu 6:} Bất phương trình nào dưới đây có nghiệm?
\begin{center}
\begin{tabular}{p{3.5cm} p{3.5cm} p{3.5cm} p{3.5cm}}
\textbf{A.} $x^2+2x+1<0$. & \textbf{B.} $x^2+2x+3\le 0$. & \textbf{C.} $x^2+3x+4<0$. & \textbf{D.} $x^2-2x\le 0$.
\end{tabular}
\end{center}
\noindent\textbf{Câu 7:} Bất phương trình nào dưới đây có nghiệm?
\begin{center}
\begin{tabular}{p{3.5cm} p{3.5cm} p{3.5cm} p{3.5cm}}
\textbf{A.} $x^2+2x+1<0$. & \textbf{B.} $x^2+2x+3\le 0$. & \textbf{C.} $x^2+3x+4<0$. & \textbf{D.} $x^2-2x+1\le 0$.
\end{tabular}
\end{center}
\noindent\textbf{Câu 8:} Bất phương trình nào dưới đây \textbf{vô nghiệm}?
\begin{center}
\begin{tabular}{p{3.5cm} p{3.5cm} p{3.5cm} p{3.5cm}}
\textbf{A.} $x^2 + 2x \ge 0$. & \textbf{B.} $x^2 + 2x + 3 \le 0$. & \textbf{C.} $x^2 - 4 < 0$. & \textbf{D.} $x^2 > 0$.
\end{tabular}
\end{center}
\noindent\textbf{Câu 9:} Bất phương trình nào sau đây \textbf{vô nghiệm}?
\begin{center}
\begin{tabular}{p{3.5cm} p{3.5cm} p{3.5cm} p{3.5cm}}
\textbf{A.} $x^2 \ge 0$. & \textbf{B.} $x^2 \le 0$. & \textbf{C.} $x^2 > 0$. & \textbf{D.} $x^2 < 0$.
\end{tabular}
\end{center}
\noindent\textbf{Câu 10:} Tập xác định của hàm số $y=\sqrt{x^2-4}$ là
\begin{center}
\begin{tabular}{p{3.5cm} p{4.5cm} p{3.5cm} p{4.5cm}}
\textbf{A.} $(-2;2)$. & \textbf{B.} $(-\infty;-2) \cup (2;+\infty)$. & \textbf{C.} $[-2;2]$. & \textbf{D.} $(-\infty;-2] \cup [2;+\infty)$.
\end{tabular}
\end{center}
\end{document}